\section{Related Work and Scope}

\paragraph{Causal structure learning from interventions.}
Observational DAG discovery is generally limited by Markov equivalence, while
known interventions can refine the equivalence class and, under additional
conditions, identify more of the graph
\citep{hauser2012interventional,yang2018interventions}.  Score-based and
continuous methods range from interventional greedy equivalence search to
differentiable optimization with interventional likelihoods
\citep{zheng2018notears,brouillard2020dcdi}.  Our task is substantially
narrower: all five coordinates and intervention targets are observed, the
candidate graphs are a finite typed motif family, and the data generator lies
in the declared transition language.  Thus the 120-world result is evidence
for recovery within that finite family, not a general DAG-discovery result and
not a claim of causal identification from observational data.

\paragraph{Causal representation learning.}
Causal representation learning asks when latent causal variables and their
relations can be recovered from high-dimensional mixtures
\citep{scholkopf2021causalrep}.  Interventional results establish
identifiability under assumptions that vary across perfect, imperfect,
unknown-target, uncoupled, and multi-node interventions
\citep{ahuja2023interventional,vonkugelgen2023nonparametric,
bing2024multinode,varici2024general}.  In contrast, this study does not infer
latent coordinates or a perceptual mixing map: the modular coordinates,
action target, and action magnitude are supplied.  Its contribution is a
controlled attribution of graph choice and typed head factorization after
those representational variables have already been declared.

\paragraph{Sparse support recovery.}
Exact support recovery by $\ell_1$ methods and greedy pursuit is known to
depend on design, signal, noise, and incoherence or restricted-curvature
conditions \citep{tropp2007omp,wainwright2009sharp,
dyer2013greedy,bahmani2013grasp}.  The seven-sparse embedding used here is an
algebraic representation over $\mathbb Z_7$, rather than a linear-Gaussian
Lasso or orthogonal-matching-pursuit problem.  Our population proposition
identifies the finite motif under full-support primitive interventions, and
its repeated-input corollary controls coordinatewise modal errors.  Neither
result proves finite random-design recovery for the implemented greedy
selector; its 120/120 recovery rate therefore remains empirical.

\paragraph{Structured world models.}
Interaction networks and graph networks encode object and relation structure
as an architectural inductive bias
\citep{battaglia2016interaction,battaglia2018relational}.  Neural relational
inference, contrastive structured world models, slot-based object discovery,
and structured exploration learn graphs or object-centric dynamics from
observations of varying complexity
\citep{kipf2018nri,kipf2020contrastive,locatello2020slot,
sancaktar2022curious}.  TSI's finite modular world is a test bed for a more
specific question: whether declared dependency information and transition-head
factorization make separately measurable contributions under controlled OOD
interventions.  The present evidence does not establish visual object
discovery, scalable graph-neural dynamics, or real-world validity.

\paragraph{Modularity and independent mechanisms.}
The independent-causal-mechanisms principle motivates decomposing a data
generating process into autonomous modules, and prior work has sought to learn
such mechanisms from transformed domains \citep{parascandolo2018independent}.
Our graph-by-head factorial design operationalizes only a restricted analogue:
it tests whether routing information still contributes when the prediction
head is correctly factorized, untied and dense, or deliberately mismatched.
Consequently, the reported interaction and stress-cohort results delimit
component attribution inside the benchmark; they do not demonstrate universal
autonomy or discovery of independent mechanisms.
